\title[Robust and Detail-Free Demand-Response Contracting]
{Distributionally Robust and Detail-Free\\ Demand-Response Contracting}



\begin{abstract}
The emission-reduction condition depends on the residual-demand distribution,
which a regulator never knows exactly. We turn this apparent weakness into two
guarantees. \emph{(R1) A worst-case certificate.} Over the natural shape class of
distributions, the \emph{uniform} is the least-favorable distribution for carbon:
it minimizes the net emission saving. Hence the single closed-form condition
$\CR\le\CR^{\dagger}$---derived for the uniform---certifies emission reduction for
the entire infinite-dimensional class, a distributionally robust guarantee.
\emph{(R2) Detail-free implementation.} Because the optimal contract's thresholds
are pinned at \emph{distribution-free} probabilities, a single quantile-targeting
mechanism implements the optimum for every distribution: the regulator needs no
knowledge of, or prior over, the demand and renewable laws. Both results are
verified numerically.
\end{abstract}



We continue the companion notes \texttt{role1\_carbon\_theory.tex} and
\texttt{role1\_contract\_theory.tex}; notation and Assumption~1 are inherited.
Two facts proved there drive everything here:
the \emph{exact} net-fossil saving
\begin{equation}\label{r1exr-eq:saving}
  \Delta(F)\eqdef\E[\text{fossil}^{\circ}]-\E[\text{fossil}^{*}]
  =\tfrac{\CM}{R-c}\big(\KFs-\KMs-\mu^{*}_{\mathrm{low}}\big)
   -\Big(\KFs-\KFo(1-\varphi^{\circ})-(\varphi^{\circ}-q^{\circ})\mu^{\circ}_{\mathrm{mid}}\Big),
\end{equation}
and the \emph{distribution-free} cutoff probabilities
$F(\KFs;\KRs)=\frac{\bp-c-\CF-\CM}{\bp-R}$,
$F(\KFs-\KMs;\KRs)=\frac{\CM}{R-c}$, $F(\KFo;\KRo)=\varphi^{\circ}$,
which hold for every $F$. In \eqref{r1exr-eq:saving} the conditional means
$\mu^{*}_{\mathrm{low}}=\E[\tilde x\mid \tilde x\in(-\KRs,\KFs-\KMs)]$ and
$\mu^{\circ}_{\mathrm{mid}}=\E[\tilde x\mid \tilde x\in(0,\KFo)]$ are the only
channels through which the \emph{shape} of $F$ enters.

%---------------------------------------------------------------------
\section{A distributionally robust emission certificate}\label{r1exr-sec:robust}

\begin{definition}[Shape ambiguity set]\label{r1exr-def:U}
Let $\Ucal$ be the set of residual-demand distributions $F$ satisfying
Assumption~1 such that the density $f(\cdot;\KRs)$ is non-increasing on the
\textsc{fcl} dispatch zone $(-\KRs,\KFs-\KMs)$ and $f(\cdot;\KRo)$ is
non-decreasing on the no-\textsc{fcl} fossil zone $(0,\KFo)$.
\end{definition}

Economically (companion Remark~10): $\Ucal$ collects grids whose deep-surplus
states are more likely than mild-surplus states, and whose high-consumption hours
are more likely than low---the empirically typical shape, with rare stress events.
The uniform sits on the boundary of $\Ucal$ (constant density).

\begin{theorem}[Least-favorable distribution; robust certificate]\label{r1exr-thm:robust}
Holding the distribution-free cutoff probabilities fixed,
\[
  \min_{F\in\Ucal}\;\Delta(F)\;=\;\Delta(F_{\mathrm{unif}}),
\]
attained by the uniform distribution. Consequently emission reduction holds for
\emph{every} $F\in\Ucal$ if and only if it holds for the uniform, i.e.\ iff
\[
  \CR\le\CR^{\dagger}
  =\tfrac12\!\Big[c+\tfrac{(R-\CM-c)^2}{R-c}+\tfrac{(R-c-\CF-\CM)^2}{\bp-R}\Big]
\]
(companion Proposition on $\CR^{\dagger}$). The single closed-form condition is a
worst-case guarantee over the infinite-dimensional class $\Ucal$.
\end{theorem}

\begin{proof}
From \eqref{r1exr-eq:saving},
$\partial\Delta/\partial\mu^{*}_{\mathrm{low}}=-\tfrac{\CM}{R-c}<0$ and
$\partial\Delta/\partial\mu^{\circ}_{\mathrm{mid}}=+(\varphi^{\circ}-q^{\circ})>0$,
the latter positive because $\varphi^{\circ}=F(\KFo;\KRo)>F(0;\KRo)=q^{\circ}$ (as
$\KFo>0$ and $F$ is increasing in $x$),
so $\Delta$ is minimized by the \emph{largest} feasible $\mu^{*}_{\mathrm{low}}$
and the \emph{smallest} feasible $\mu^{\circ}_{\mathrm{mid}}$. A non-increasing
density on an interval has conditional mean at most the interval midpoint, with
equality iff the density is constant; a non-decreasing density has conditional
mean at least the midpoint, with equality iff constant. Hence over $\Ucal$,
$\sup\mu^{*}_{\mathrm{low}}$ and $\inf\mu^{\circ}_{\mathrm{mid}}$ are \emph{both}
the respective midpoints, attained \emph{simultaneously} by the uniform. So the
uniform jointly minimizes $\Delta$. The closed-form $\CR^{\dagger}$ is the uniform
threshold; if it certifies $\Delta\ge0$ at the worst case, $\Delta(F)\ge0$ for all
$F\in\Ucal$.
\end{proof}

\begin{remark}[Why this is the right robustness statement]
Without a shape restriction the ambiguity is vacuous---an adversary placing all
conditional mass at an interval endpoint can force $\Delta<0$, so no
distribution-free certificate can exist (verified: unrestricted worst case is
negative). $\Ucal$ is the minimal economically meaningful class that restores a
certificate, and within it the certificate is \emph{exactly} the uniform
closed form. Theorem~\ref{r1exr-thm:robust} thus identifies the uniform benchmark of the
original paper not as a convenient special case but as the \emph{binding} worst
case for a whole family---the strongest reading the closed form can carry.
\end{remark}

\begin{remark}[Stochastic-order reading]
Because $\Delta$ moves with the two conditional means, and conditional means on a
fixed interval are ordered by the dispersive/likelihood-ratio order, $\Delta$ is
monotone along any one-parameter tilt of the density toward or away from the
uniform. The uniform is the fixed point of this monotone family---the least
disperse member consistent with $\Ucal$'s monotonicity.
\end{remark}

%---------------------------------------------------------------------
\section{Detail-free implementation}\label{r1exr-sec:detailfree}

A regulator does not know $F$. Remarkably, it need not.

\begin{theorem}[Prior-independent optimal mechanism]\label{r1exr-thm:detailfree}
Define the \emph{quantile-targeting mechanism}: choose capacities so that, at the
realized residual-demand law,
\[
  \Pr(\tilde x\le \KF)=\frac{\bp-c-\CF-\CM}{\bp-R},\quad
  \Pr(\tilde x\le \KF-\KM)=\frac{\CM}{R-c},\quad
  \Pr(\tilde x\le \KFo)=\varphi^{\circ},
\]
and dispatch $\dstar(x)=\min\{(\KF-x)_+,\KM\}$. For \emph{every} distribution $F$
satisfying Assumption~1, this mechanism implements the optimal contract and the
optimal capacities. It uses no knowledge of, and no prior over, $F$: only the
right-hand-side probabilities, which are functions of the primitives
$(R,c,\CF,\CM,\beta,p)$ alone.
\end{theorem}

\begin{proof}
The capacity first-order conditions (companion Lemma) equate $F$ at each threshold
to the stated probability; these probabilities contain no distributional object.
Targeting them therefore selects the optimal $(\KRs,\KFs,\KMs)$ whatever $F$ is,
and Theorem~1 of the contract note makes $\dstar$ optimal given those capacities.
A regulator can hit the targets from observed empirical quantiles of residual
demand without estimating $F$.
\end{proof}

\begin{remark}[Wilson doctrine]
Theorem~\ref{r1exr-thm:detailfree} is a detail-free implementation in the sense of the
Wilson doctrine: the optimal mechanism is \emph{prior-independent}. The economic
content of the distribution-free cutoffs (companion Lemma) is precisely this---the
margins $R-c$, $\bp-R$, and the no-\textsc{fcl} cutoff probability $\varphi^{\circ}$
summarize everything
the regulator must know. Distributional detail matters for the \emph{magnitude} of
the carbon saving (Theorem~\ref{r1exr-thm:robust}), never for \emph{how to act}.
\end{remark}

\begin{corollary}[Robust+detail-free package]
Combining Theorems~\ref{r1exr-thm:robust}--\ref{r1exr-thm:detailfree}: the regulator implements
the utility's \emph{profit-maximizing} demand-response contract knowing only six
scalar primitives, and---provided the grid lies in the shape class $\Ucal$ and
$\CR\le\CR^{\dagger}$---this privately optimal contract is \emph{guaranteed} to
lower emissions regardless of the true demand and renewable laws. Emission
reduction is a consequence of the profit-optimal contract under the stated
condition, not a separate regulatory objective.
\end{corollary}

%---------------------------------------------------------------------
\section{Numerical validation}\label{r1exr-sec:num}

\texttt{verify\_robust.py}:
\begin{itemize}[leftmargin=1.6em]
  \item Over $3000$ random non-increasing densities on the \textsc{fcl} zone, the
        conditional mean is $\le$ midpoint in $3000/3000$ cases, with the supremum
        equal to the midpoint (uniform extremal); symmetrically, over $3000$ random
        non-decreasing densities the conditional mean is $\ge$ midpoint in
        $3000/3000$, infimum at the midpoint.
  \item Tilting the density continuously from uniform (shape $0$) into $\Ucal$, the
        net carbon saving $\Delta$ is \emph{minimized at the uniform}---confirming
        the uniform is least-favorable (Theorem~\ref{r1exr-thm:robust}).
  \item The cutoff probabilities of Theorem~\ref{r1exr-thm:detailfree} are the
        distribution-free values verified to $10^{-16}$ on the non-uniform
        $\mathrm{Beta}(2,2)$ scale-mixture in the companion \texttt{verify\_exact.py}.
\end{itemize}

\paragraph{Scope.} Strictly A1 (general contracting theory and numerical methods):
the object is the utility's profit-maximizing contract under \emph{distributional}
ambiguity (ambiguity over the demand/renewable law $F$), and its prior-independent
implementation. This is distinct from---and does not enter---A2 (empirical
calibration / magnitudes), B1 (reliability, welfare, and the social-vs-private
optimum, including corrective pricing), or B2 (the flexibility continuum, storage,
and the multi-grid / price-exogeneity / dynamics robustness-and-limitations
discussion). The robustness here is robustness of the \emph{contract} to the
unknown distribution, not B2's robustness of the model to market structure.
